% ===========================================================================
% Motivating Example -- paper-ready LaTeX
% Pairs with Figure 1 (motivation_figure.tex, Skill `lean4-proof`).
% Replace \ref{sec:corpus}, \ref{sec:taxonomy}, \ref{fig:motivation}.
% No extra packages beyond Figure 1's tikz/xcolor; text needs only the kernel.
% ===========================================================================
\section{Motivating Example}
\label{sec:motivation}
%
% Local only: a long unbreakable \texttt{} token at the end of a narrow measure
% would otherwise overflow. Verified clean from 8cm to 16cm text width.
\begingroup\emergencystretch=3em

An Agent Skill is a directory an agent loads as instructions: a \texttt{SKILL.md}
contract, together with the scripts and reference documents that contract points
to. The agent cannot execute the contract to see whether it holds. It reads the
text and acts on it. We take one Skill from the SkillBench corpus we audit in
Section~\ref{sec:corpus} and read its contract the way an agent would.

\texttt{lean4-proof} helps an agent write Lean~4 proofs. It is 28 files, 26 of
them Markdown. Line~18 of the proving Skill's \texttt{SKILL.md} describes an ``Interactive
Commands'' table of ``10 slash commands''. Line~32 of the same file advertises
``7 slash commands''. Fourteen lines apart, the document gives
two answers, and nothing else in the Skill says which one holds.

A second inconsistency sits between the documentation and the script it ships.
the memory Skill's \texttt{SKILL.md:10} states that it ``enables persistent
learning and knowledge accumulation across Lean 4 formalization sessions''. The
helper that is supposed to provide this has its MCP call commented out at
line~90 and prints, at line~94, \texttt{Pattern stored successfully
(simulation)}. The agent that runs the helper is told the write
succeeded. Nothing is written, and nothing fails.

The third runs against a service the Skill does not own. The pattern reference
calls \texttt{memory.\allowbreak{}search\_\allowbreak{}entities(\ldots)}; the Skill configures the
community MCP server \texttt{@modelcontext\allowbreak{}protocol/\allowbreak{}server-\allowbreak{}memory}
(\texttt{SKILL.md:371}), whose tools carry other names. The documented
interface is not the interface the pinned service provides.

All three are settled by reading the files, none of them raises an error when
the Skill is used, and each one breaks a different pair of artifacts: the
document against itself, the document against its code, and the Skill against
its environment. All three are instances of the types we define in
Section~\ref{sec:taxonomy} and carry high severity
(Figure~\ref{fig:motivation}).

\paragraph{Observation 1.} \emph{A Skill contract can contradict itself.} The
command count above is stated twice, differently, in a single file. No
documentation tool we know of looks for this: link checkers, linters and doc
builds all pass. If a Skill is treated as one trusted specification, then a
contract that disagrees with itself has no correct reading, and the agent cannot
tell which half to follow. RQ1 asks which kinds of contract inconsistency occur
in real Skills, how often, and how they distribute across the three boundaries.

\paragraph{Observation 2.} \emph{The damaging ones are self-certifying.} The
memory helper does not merely fail to persist: it prints a success message and
exits cleanly. An agent has no signal to act on, and neither does the caller that
reads its output. Detection therefore cannot stop at finding a mismatch; it has
to separate the ones an agent will build on from the ones it will read past. RQ2
asks whether these inconsistencies can be found automatically, and whether that
ranking can be recovered from the artifacts.

\paragraph{Observation 3.} \emph{The cost is deferred and the repair is not
obvious.} An agent that trusts the memory Skill will store nothing, believe it
has stored something, and lose the knowledge when the session ends. The failure
surfaces later, in a different session, with no path back to its cause. Repair
requires deciding which side is authoritative, and for the third defect that
side is a third-party server the Skill does not control. RQ3 asks what an
inconsistency costs an agent, and whether repairing it recovers that cost.
\endgroup
